\section{Coset geometries and the objects retained by an
operation}\label{coset-geometries-and-the-objects-retained-by-an-operation}

A geometry has a finite set of types; its flags are sets of pairwise
incident objects of distinct types, every flag extends to a chamber, and
its residues are the induced geometries on objects incident with a fixed
flag. It is thin when every panel has exactly two completions and
residually connected when every residue of rank at least two has
connected incidence graph. Regularity means chamber transitivity of the
full type-preserving automorphism group. No finiteness of the object set
is implicit.

A marking of a group G in the geometric sections is a finite indexed generating family of nonidentity
involutions \((r_i)_{i\in I}\). Write \(G_J=\langle r_j:j\in J\rangle\).
The intersection property is \[
G_J\cap G_K=G_{J\cap K}\qquad(J,K\subseteq I).
\tag{C1}
\] The maximal parabolics are \(P_i=G_{I\setminus\{i\}}\); their right
cosets are incident when they intersect. Typed coset Helly means that
every pairwise intersecting family with distinct types has a common
element. UC.Geometry shows that IP and typed Helly give exactly the
regular thin residually connected coset geometry, with G as its full
type-preserving group. IP alone is not sufficient for that conclusion.

\textbf{Proposition UC.Geometry (coset recognition and faithful chamber
action).} Suppose the marking of \(G\) consists of nonidentity
involutions, has IP, and has typed coset Helly for its maximal
parabolics. The coset incidence system is then a thin, residually
connected geometry, and \(G\) is its full type-preserving automorphism
group, acting regularly on chambers. More generally, the full
type-preserving automorphism group of any thin, residually connected
geometry of finite rank acts freely on chambers. A chamber-transitive
subgroup is consequently the full group.

\emph{Proof.} In the coset system, a common coset point extends every
flag to a chamber by choosing the coset of each missing type through
that point. Right translation sends every such flag to its base flag.
For a base flag of type set \(S\), IP gives stabilizer
\(G_{I\setminus S}\). In particular the full chamber stabilizer is
\(1\), while a panel stabilizer is a group generated by one nonidentity
involution and has order two. There are therefore exactly two chambers
through each panel.

Consider a residue with remaining type set \(T\), where \(|T|\geq2\).
Its stabilizer is \(G_T\), and its maximal object stabilizers are
\(G_{T\setminus\{i\}}\), \(i\in T\). Every marker indexed by \(T\) lies
in at least one of these proper parabolics. A word in those subgroups
supplies a walk between their translated base objects in the residue
incidence graph; all objects are translates of base objects and all base
objects lie in a common flag. The residue is therefore connected.

For the general assertion, first observe that every residually connected
finite-rank geometry has a connected chamber graph. This follows by
induction on rank: chambers through a fixed object are connected through
galleries in its residue, and a path in the incidence graph, extending
each consecutive incident pair to a chamber, connects arbitrary chambers
through these subgraphs. The rank-one case holds because all chambers
share the empty panel. In a thin geometry an automorphism fixing a
chamber fixes its unique neighbor across each panel. It therefore fixes
every chamber along every gallery, and hence every object. The chamber
action is free. If a subgroup is chamber-transitive, any full
automorphism agrees on a base chamber with one subgroup element; their
quotient fixes that chamber and is the identity. Apply this to the coset
geometry to obtain the stated full automorphism group. \(\square\)

\textbf{Lemma C.Gallery (the marked stabilizers).} In a thin regular
residually connected geometry, the full group acts freely on chambers.
At a base chamber let \(r_i\) be the unique group element carrying it
across its i-panel. Then \(r_i\) is an involution, the \(r_i\) generate
the full group, and the stabilizer of a base subflag F is \[
G_F=\langle r_i:i\notin t(F)\rangle.
\tag{C2}
\]

\emph{Proof.} UC.Geometry proves chamber-graph connectedness and
freeness of the chamber action. The square of \(r_i\) fixes the base
chamber, so \(r_i\) has order two. A chamber gallery expresses every
chamber image by successive panel changes, which gives generation. Apply
the same gallery argument in the residue at F. Only the missing types
can change, so its chamber-transitive stabilizer is generated by their
\(r_i\). Conversely those \(r_i\) fix F. This proves (C2). \(\square\)

\textbf{Lemma UC.Helly (anchored criterion).} A finite indexed subgroup
family \((P_i)_{i\in I}\) has typed coset Helly if and only if the
following condition holds. For every \(i\in I\), every
\(J\subseteq I\setminus\{i\}\) with \(|J|\geq2\), and any selected
\(j_0\in J\), each \(g\in P_{j_0}\) satisfying \(g\in P_iP_j\) for all
\(j\in J\) admits \(k\in\bigcap_{j\in J}P_j\) with \(P_i g=P_i k\).

\emph{Proof.} The cosets \(P_i g\) and \(P_j\) intersect precisely when
\(g\in P_iP_j\). Thus Helly gives a common element \(k\), proving
necessity. Conversely, induct on the number of cosets in a pairwise
intersecting family. A family of at most two already has a common
element. For a larger family, the induction hypothesis gives a common
element of all but one coset. Right multiplication by its inverse turns
those cosets into the base cosets \(P_j\). The remaining coset meets
\(P_{j_0}\), so it has a representative \(g\in P_{j_0}\). Its
intersections with the other base cosets are precisely the product
conditions in the statement. The criterion supplies \(k\), and undoing
the right translation gives a common element of the original family.
\(\square\)

\textbf{Lemma C.HellyCount (finite recognition).} For a finite group,
put \[
X_{ij}=\{P_i x:x\in P_j\},\qquad P_J=\bigcap_{j\in J}P_j.
\] Typed coset Helly is equivalent to \[
\left|\bigcap_{j\in J}X_{ij}\right|=[P_J:P_J\cap P_i]
\quad(i\notin J,\ |J|\ge2).
\tag{C3}
\]

\emph{Proof.} The cosets represented in \(P_J\) form a subset of the
displayed intersection, of size the right side. Equality is therefore
equality of these sets. This is exactly the anchored condition of
UC.Helly: a coset meeting every base \(P_j\) has a representative in
their common subgroup. The finite count proves both directions.
\(\square\)

\textbf{Theorem UC.Action (faithful action after retaining objects).}
Let \(G\) act on sets \(X\) and \(Y\). Suppose a \(G\)-invariant subset
\(Y_0\subseteq Y\) has an equivariant surjection onto \(X\). Suppose
also that \(Y\) has an injective equivariant encoding in a tagged
disjoint union of finite products, each factor of which is \(X\) or its
power set, and that the tags are fixed by \(G\). Then \[
\ker(G\curvearrowright Y)=\ker(G\curvearrowright X)
=\bigcap_{x\in X}G_x.
\tag{5.1}
\] If \(X\) is a disjoint union of transitive orbits with selected
stabilizers \(H_j\), the common kernel is \[
N=\bigcap_j\operatorname{core}_G(H_j)
=\operatorname{core}_G\left(\bigcap_jH_j\right).
\tag{5.2}
\] The faithful permutation image is \(G/N\). If its action on \(Y\) is
chamber-transitive on a thin, residually connected finite-rank geometry,
then \(G/N\) is the full type-preserving automorphism group. The
preimages in \(G\) of a chamber stabilizer and of a panel stabilizer are
respectively \(N\) and a subgroup of index two over \(N\).

\emph{Proof.} An element fixing \(Y\) pointwise fixes \(X\) through the
equivariant surjection. Conversely, an element fixing \(X\) fixes all
subsets of \(X\), every finite tuple formed from these objects, and all
fixed tags. Injectivity of the encoding then forces it to fix \(Y\)
pointwise. This proves (5.1). In a transitive orbit all point
stabilizers are conjugates of the selected one, so intersection over all
points gives the first equality in (5.2). The second is the intersection
of the same conjugate subgroups in the opposite order. The first
isomorphism theorem identifies the faithful image.

By UC.Geometry that image has trivial chamber stabilizers and is the
full type-preserving group. Its panel stabilizer acts transitively on
the two chambers through the panel, with trivial stabilizer of either
chamber, so has order two. Taking preimages proves the final assertion.
\(\square\)

Independently labeled copies of the same subset need extra data in their
encoding; they cannot be silently treated as the subset itself.

\textbf{Proposition UC.ExtensionOrbits (classification above a retained
base).} Let a group \(B\) act on sets \(Z,E\), and let \(p:Z\to E\) be
equivariant. For one representative \(e\) of each orbit in
\(B\backslash E\), there is a natural bijection \[
B\backslash Z\ \cong\ \coprod_{e\in B\backslash E}B_e\backslash p^{-1}(e).
\tag{5.6}
\]

\emph{Proof.} Translate any \(z\in Z\) so that \(p(z)\) is the selected
representative of its orbit. If two translated points over that
representative are \(B\)-equivalent, an element taking one to the other
fixes their image \(e\), so belongs to \(B_e\). Conversely a
\(B_e\)-equivalence is already a \(B\)-equivalence. These two
observations prove surjectivity and injectivity of (5.6). \(\square\)
